\documentclass[11pt]{article}
\usepackage[utf8]{inputenc}
\usepackage[T1]{fontenc}
\usepackage{amsmath,amssymb,amsthm}
\usepackage{booktabs,tabularx}
\usepackage{graphicx}
\usepackage{xcolor}
\usepackage{xurl}
\usepackage{hyperref}
\usepackage[capitalise,noabbrev]{cleveref}
\usepackage{enumitem}
\usepackage{microtype}
\usepackage{natbib}
\usepackage{authblk}
\usepackage[a4paper,margin=1in]{geometry}

\hypersetup{
  colorlinks=true,
  linkcolor=blue!70!black,
  citecolor=blue!70!black,
  urlcolor=blue!70!black
}

\setlist{nosep}
\newcommand{\R}{\mathbb{R}}
\newcommand{\N}{\mathbb{N}}
\newcommand{\bx}{\mathbf{x}}
\newcommand{\bh}{\mathbf{h}}
\newcommand{\bg}{\mathbf{g}}
\newcommand{\bW}{\mathbf{W}}
\newcommand{\cL}{\mathcal{L}}
\newcommand{\supp}{\operatorname{supp}}

\newtheorem{definition}{Definition}
\newtheorem{theorem}{Theorem}
\newtheorem{proposition}{Proposition}

\date{August 2026}


\hypersetup{
  pdftitle={Binary Activation: Identity-Bound FFN Confinement for Shared Neural Networks},
  pdfauthor={Ryan McCormick},
  pdfsubject={Credential-derived binary activation and FFN confinement}
}

\title{\textbf{Binary Activation}\\
\large Identity-Bound FFN Confinement for Shared Neural Networks}
\author[1]{Ryan McCormick}
\affil[1]{Independent Researcher}

\begin{document}
\maketitle

\begin{abstract}
Public-key infrastructure made identity administrable by binding keys to names
across organizations. Shared model artifacts need an analogous addressable
structure for capability: a signed grant must identify not only who may execute,
but which internal support it authorizes. Binary activation supplies that
structural basis at declared neural boundaries.

Shared neural networks need a mechanism that determines not only who may call a
model, but which learned coordinates a credential may activate. We call this
mechanism \emph{binary activation}. An authenticated credential resolves to a
binary mask over a feed-forward network (FFN) activation. The model computes
$\tilde{\bh}=\bh\odot\bg_r$: coordinates assigned to regime $r$ are preserved,
and every excluded coordinate is exactly zero. The same multiplication governs
backpropagation. When the gate is placed on the expanded FFN activation before
the down projection, loss gradients and optimizer state can be confined to the
aligned columns of the first projection, hidden bias entries, and rows of the
second projection.

We formalize the local gate, gradient, and aligned-weight statements in Lean 4
and report seven complementary experiment families. A five-seed formative study
finds post-encoder co-training costs of 0.14--0.38 percentage points across
$R=8$--128; no paired test survives Bonferroni correction. In the stricter
frozen-backbone DistilBERT study, every intermediate
FFN is gated and only authorized FFN slices plus a private classifier are
trainable. Across three seeds, owning accuracy falls from $91.28\%$ at $R=1$ to
$86.21\%$ at $R=16$, while frozen shared weights, off-partition parameters and
optimizer moments, and inactive classifiers remain bit-exact unchanged. Private
residual adapters recover $90.47\%$ owning accuracy at $R=4$ without modifying
inactive lanes. A capacity-preserving synthetic classifier with $d=10R$ obtains
256/256 owning predictions at $R=128$ and zero hits in 32,512 wrong-key probes;
authorized support extraction has maximum logit difference zero. It is a
memorization and scaling result, not held-out utility evidence. Exact
extraction, full-capacity inference controls, and
deliberately broken placement controls distinguish zero-cost authorization and
reparameterization from the measured cost of dividing FFN capacity. The claim
is intentionally local: binary activation supplies exact coordinate exclusion
at a correctly placed FFN gate, not complete-model privacy through shared
attention, caches, logs, or undeclared bypasses.
\end{abstract}

\section{Introduction}

Public-key infrastructure answers whose key is being presented and who vouches
for it. As trained models are licensed, co-located, and embedded in systems
outside their creator's custody, an analogous question arises: which capability
inside this artifact may this authenticated principal execute? Encryption and
endpoint admission alone do not provide an internal object to which a scoped,
signed grant can refer. Binary activation supplies such an addressable object at
a declared model or FFN boundary.

Multi-tenant inference commonly treats authorization as a routing decision:
after a caller is admitted, a serving layer chooses a model, adapter, or data
source. That is useful, but it does not define a structural boundary inside a
shared learned representation. Binary activation connects an authenticated
principal to a computation-level decision. A credential is verified by a
trusted authority, resolved to a regime, and converted into a binary support
mask. The model can then use an FFN coordinate if and only if the resolved mask
activates it.

The core operation is elementary:
\begin{equation}
  \tilde{\bh}=\bh\odot\bg_r,
  \qquad \bg_r\in\{0,1\}^d.
  \label{eq:gate}
\end{equation}
Its value comes from where it is placed and how it is controlled. At the gate,
$0\cdot x=0$ eliminates an excluded coordinate for every input value, and
$1\cdot x=x$ preserves an active coordinate. By the chain rule, the same mask
zeros the corresponding loss gradient. With an aligned optimizer, one
credential can therefore read from and write to only its declared FFN support.

This paper contributes:
\begin{itemize}
  \item a capability-PKI formulation in which signed, identity-bound grants
  address declared activation supports inside a model artifact;
  \item a deterministic keyed construction of disjoint and exhaustive regime
  masks;
  \item a precise FFN gate placement with forward, gradient, and aligned-state
  boundaries;
  \item Lean 4 formalization of the central algebraic statements;
  \item formative and strict DistilBERT classification studies measuring
  utility from $R=1$ through $R=128$ at two declared gate surfaces;
  \item a capacity-preserving classifier construction that holds each regime's
  allocated width fixed while total width grows as $d=qR$;
  \item private-capacity lanes and exact inactive-state checks;
  \item exact extraction and full-capacity inference controls; and
  \item deliberately broken placement and bypass controls.
\end{itemize}

The central security statement is structural but narrow. It applies at the
declared FFN activation and aligned trainable state. Identity verification,
credential custody, mask dispatch, optimizer behavior, gate placement, and
closure of bypass paths remain part of the trusted implementation boundary.

\section{Classifier and Modular Use Cases}

\subsection{Why the DistilBERT experiments worked}

DistilBERT is a Transformer. The relevant boundary is therefore not
``Transformer'' versus ``non-Transformer,'' but which part of the computation
the gate governs. In the formative classifier study, DistilBERT supplies a
shared Transformer encoder and the gate acts on the final CLS representation.
In the strict study, a gate acts inside each Transformer block's FFN while
attention, normalization, embeddings, and residual paths are frozen and
shared. Neither construction claims to partition attention.

The formative post-encoder result worked because the complete encoder was
co-trained with the masked CLS representation. Optimization could therefore
place discriminative features onto the supports that each regime would retain.
That experiment measures utility under a learned support constraint; because
all regimes update the shared encoder, it is not evidence of separated tenant
state inside DistilBERT.

The strict result works for a different reason. A Transformer FFN already
constructs an explicit expanded tensor after an element-wise nonlinearity and
before a linear down projection. Its coordinates, aligned projection slices,
and optimizer entries are identifiable constituent pieces on which a binary
support can act. Freezing attention, embeddings, normalization, and residual
parameters leaves a shared feature-producing backbone while the governed FFN
slices and private classifier adapt. The resulting capacity curve is therefore
a valid Transformer-classification result, not an attention-partition claim.

\subsection{Classifier uses and capacity policies}

This still covers a broad and commercially useful classifier surface. Examples
include identity-scoped intent and document routing, fraud and abuse scoring,
policy and entitlement decisions, risk tiers, content moderation, and
classification heads attached to vision or audio encoders. A regime need not
mean one entire tenant: the authority may bind a support to a tenant plus a
sub-identity, role, licensed capability, or data purpose, and may authorize a
declared union when one principal is entitled to several capabilities.

Three patterns avoid treating a model as one indivisible dense object:
\begin{itemize}
  \item an atomic whole-model gate authorizes an unchanged classifier before
  execution and therefore does not alter its numerical computation;
  \item a shared frozen encoder feeds a gated classifier, FFN slice, private
  adapter, or private expert; and
  \item a block-sparse classifier or expert bank is sized with a fixed width
  $q$ per regime, so total governed width is $d=qR$ rather than a fixed $d$
  divided into progressively smaller pieces.
\end{itemize}
The third pattern preserves allocated per-regime capacity by spending storage
linearly in $R$. It is natural for branched classifiers, expert banks, adapters,
and other modular surfaces. Binary activation supplies identity-bound selection
and exact zeroing at the declared gate; sparse dispatch or authorized extraction
may avoid evaluating inactive blocks. It does not make additional private
capacity free.

\subsection{Why attention is a different space}

Binary Activation requires more than a tensor on which multiplication is
defined. The protected computation must expose stable constituent pieces: an
indexed support before cross-support mixing, aligned parameters and optimizer
state, and no ungated path around the support boundary. Standard self-attention
does not expose those pieces as fixed tenant coordinates. Queries and keys form
input-dependent pairwise dot products. Those products encode both vector norms
and angular proximity; when norms are controlled, the angular term corresponds
to cosine adjacency. Softmax then
couples every candidate through a shared denominator, the output projection
recombines heads, and normalization and residual additions mix or bypass the
result.

Masking an arbitrary attention coordinate still produces an exact zero at that
coordinate, but it does not construct an independent tenant constituent from a
distributed relational representation. A fixed mask can change token
affinities and the softmax distribution globally, which explains why post-hoc
head ablation damages utility. To make attention partitionable, the architecture
would have to construct complete lanes from the outset---including Q/K/V,
output projection, normalization, FFN, residual stream, and cache state---and
gate every route before recombination. That is possible as a research design,
but it is a different and substantially higher-risk construction than the
DistilBERT classifier experiments reported here.

\section{Binary Activation Model}

\subsection{Credential-to-support resolution}

Let $c$ be an authenticated credential associated with a user, agent, or
tenant. A trusted authority resolves $c$ to a regime identifier $r$ and the
corresponding mask $\bg_r$. Callers do not supply masks or choose the underlying
coordinate assignment. The model input is therefore the pair $(\bx,c)$, not
$\bx$ alone:
\begin{equation}
  c\longmapsto r\longmapsto\bg_r,
  \qquad M(\bx,c)=M(\bx;\bg_r).
\end{equation}

Authentication and algebra perform different jobs. Authentication determines
which support is authorized; multiplication enforces that support at the
declared tensor. Knowledge of mask indices is not itself a credential when the
runtime refuses caller-supplied assignments.

The resolved regime then becomes an authenticated execution identity for each
downstream gate that verifies the same bound grant, model, operation, and policy
scope. The resulting chain is external authentication, boundary authorization,
and downstream regime authentication. A bare regime integer or activation
tensor does not authenticate itself. Identity assurance terminates at the
verifier-configured PKI root and the operator's certificate and key-custody
controls; PKCS\#12 credential packaging alone does not prove hardware residency.

\subsection{Disjoint keyed assignment}

Given hidden dimension $d\ge R$, a gate partition is a family
$\{\bg_0,\ldots,\bg_{R-1}\}\subset\{0,1\}^d$ satisfying
\begin{equation}
  \bg_r\odot\bg_s=\mathbf 0\quad(r\ne s),
  \qquad \sum_{r=0}^{R-1}\bg_r=\mathbf 1.
  \label{eq:partition}
\end{equation}
Every coordinate belongs to exactly one regime. Equal-width partitions require
$R\mid d$, equivalently $d=qR$ for an integer $q\ge1$; otherwise block sizes
may differ by at most one.

The reference construction uses HMAC-SHA256 as a pseudorandom function to drive
Fisher--Yates over the coordinate indices. The encoded context includes a model
digest, model version, $R$, $d$, the shuffle position, and a rejection-sampling
counter. Accepted 256-bit words are reduced modulo the current shuffle range
only after rejecting the incomplete residue tail. The resulting permutation is
partitioned into $R$ blocks and converted to binary masks.

The keyed shuffle makes assignment deterministic for the same key and context
and unpredictable under the standard PRF assumption. Disjointness and
exhaustiveness come from partitioning the permutation; the local zero at the
gate comes from the binary mask, not from secrecy of the permutation.

\subsection{Computational overhead}

The runtime gate is one element-wise multiply over $d$ coordinates per gated
tensor: $d$ scalar multiplications with no accumulation, rather than an extra
matrix product. Writing the operation as
$\operatorname{diag}(\bg_r)\bh$ and evaluating that dense matrix would cost
$\Theta(d^2)$ operations and would be an implementation error. The intended
broadcast multiply is $\Theta(d)$ against the
$\Theta(d\,d_{\mathrm{in}})$ and $\Theta(d\,d_{\mathrm{out}})$ projection
costs around it, and it may be fused with a surrounding kernel.

For inference-output equivalence, a constant mask may instead be folded into
the aligned rows of the down projection. That packaging transformation removes
the per-request multiply, but it does not instantiate the declared
activation-level zero or the training-state confinement boundary; it is an
extracted or reparameterized artifact. A bit-packed mask requires $d$ bits
(96 bytes at $d=768$), while an expanded Boolean or floating-point tensor uses
more implementation-dependent storage. No dedicated gate-overhead
microbenchmark is included, so this paper makes no numerical latency claim.

\section{FFN Placement and Confinement}

\subsection{Intermediate-FFN activation}

For row-vector activations, write a two-projection FFN as
\begin{align}
  \bh &= \sigma(\bx\bW_1+\mathbf b_1), \\
  \tilde{\bh} &= \bh\odot\bg_r, \\
  \hat{\mathbf y} &= \tilde{\bh}\bW_2+\mathbf b_2.
  \label{eq:ffn}
\end{align}
The gate lies after the element-wise nonlinearity and immediately before the
down/output projection. An excluded hidden coordinate contributes zero to every
output coordinate regardless of its pre-gate value.

The location is essential. If normalization or another mixing operation moves
information between coordinates before the mask is applied, a later zero cannot
undo signal already transferred into active coordinates. Likewise, an ungated
residual, adapter, cache, or other parallel branch lies outside the claim.

\subsection{Forward and gradient properties}

For coordinate $j$:
\begin{align}
g_{r,j}=0 &\implies (\bh\odot\bg_r)_j=0, \\
g_{r,j}=1 &\implies (\bh\odot\bg_r)_j=h_j.
\end{align}
Backpropagation gives
\begin{equation}
  \frac{\partial\cL}{\partial\bh}
  =\frac{\partial\cL}{\partial\tilde{\bh}}\odot\bg_r,
  \label{eq:gradient}
\end{equation}
so $g_{r,j}=0$ implies $\partial\cL/\partial h_j=0$ for every loss,
input, and batch size.

In the notation of \cref{eq:ffn}, the aligned loss gradients are therefore zero
for excluded columns $\bW_1[:,j]$, hidden-bias entry $b_{1,j}$, and rows
$\bW_2[j,:]$. The output bias $\mathbf b_2$ is not indexed by the hidden
partition and is shared unless separately frozen, partitioned, or gated.

\subsection{Parameter state and optimizer state}

Zero loss gradient alone does not stop every optimizer from moving a parameter.
Decoupled weight decay, momentum accumulated under another regime, or an
unmasked post-processing step can modify an inactive coordinate. Exact
parameter-state confinement therefore requires one of the following equivalent
implementation patterns:
\begin{itemize}
  \item optimizer moments stored and updated only for the active support;
  \item inactive weights and moments restored after each step; or
  \item an optimizer whose complete update, including weight decay, is masked
  by the same regime support.
\end{itemize}

Under that condition, regime-$r$ training updates only the aligned FFN slices
and any explicitly private state such as a regime-selected classifier. Shared
attention, embeddings, normalization, output heads, and runtime state do not
inherit this guarantee.

\subsection{Per-sample co-training}

Binary activation supports mixed-regime minibatches. For examples
$\{(\bx_i,y_i,r_i)\}_{i=1}^B$, each sample uses its own mask:
\begin{equation}
  \tilde{\bh}_i=\bh_i\odot\bg_{r_i}.
\end{equation}
The resulting gradients can be accumulated in one optimizer step while each
sample remains confined to its support. The supports may be assigned to
different tenants or to multiple capabilities held by one principal. A runtime
that intentionally authorizes a union removes the boundary among the included
supports; coordinates outside the union remain excluded.

\section{Formal Verification}

The accompanying pinned Lean 4 project builds 21 modules and 312 active theorem or
lemma declarations, with no active \texttt{sorry} or \texttt{admit}
declarations. A checked-in \texttt{\#print axioms} audit shows that the nine
headline gate and placement statements use only Lean/Mathlib's standard
logical axioms. Five project-specific \texttt{axiom} declarations are present
in separate conditional or historical theorem chains; none reaches the
headline statements used here.

The selected results are:
\begin{itemize}
  \item \texttt{gradient\_isolation}: an inactive gate coordinate has zero
  activation gradient;
  \item \texttt{weight\_update\_confined}: first-projection updates are
  confined to active hidden columns;
  \item \texttt{w2\_update\_confined}: second-projection updates are confined
  to active hidden rows;
  \item \texttt{masks\_orthogonal}: distinct regime masks have disjoint
  support;
  \item \texttt{residual\_block\_full\_confinement}: the aligned FFN residual
  branch remains confined for arbitrary upstream gradient;
  \item \texttt{non\_preserving\_breaks\_w1\_confinement}: an intervening
  non-support-preserving operation such as LayerNorm invalidates the first
  projection confinement statement; and
  \item \texttt{compressed\_eq\_gated}: selecting active support yields the
  same idealized linear result as full gated execution.
\end{itemize}

These are local algebraic theorems. They do not prove the correctness of an
identity provider, key service, tensor-axis mapping, optimizer implementation,
or the absence of bypass paths.

\section{Experimental Evaluation}

\subsection{Evidence map}

The experiments answer different questions and are not pooled into one claim.
The first study measures whether gate-aware optimization can use very narrow
supports. The wide-classifier study instead holds each regime's allocated width
fixed while total width grows with $R$. The strict study then moves the gate
inside every FFN and audits
parameter and optimizer state. Private lanes test the effect of restoring
per-regime capacity. Extraction and full-capacity controls test exact inference
preservation when the authorized computation is not capacity constrained.
Finally, negative controls show that the harness detects invalid placement and
an ungated residual bypass.

\begin{table}[ht]
\centering
\caption{Completed experiment families included in this evaluation.}
\label{tab:evidence-map}
\small
\begin{tabularx}{\textwidth}{@{}p{0.34\textwidth}X@{}}
\toprule
Experiment and repetition & Headline result \\
\midrule
Post-encoder co-training: 5 paired seeds; 7 values of $R$ &
Observed gap 0.143--0.376 percentage points \\
Capacity-preserving classifier: 1 deterministic $R=128$ run &
256/256 owning; 0/32,512 wrong-key hits; selected/dense difference 0 \\
Strict intermediate FFN: 3 seeds; $R=1,2,4,8,16$ &
Bit-exact inactive state; measured capacity curve \\
Private residual lanes: 3 seeds; 2 designs at $R=4$ &
$90.47\%$ adapter accuracy; inactive-lane delta 0 \\
Authorized extraction: 10 seeds; 160 comparisons &
80/80 fp32 bit-identical; 80/80 fp16 within tolerance \\
Full-capacity controls: 1 trained model plus local authorization probe &
Zero measured inference difference when capacity is retained \\
Placement controls: 1 deterministic local protocol run &
Correct gate 0; invalid LayerNorm and residual paths nonzero \\
\bottomrule
\end{tabularx}
\end{table}

All reported values are backed by checked-in runners and machine-readable
artifacts. The formative statistics are in
\path{experiments/results/series1_results_20260803_202919.json}; strict FFN and
private-lane runs are in the corresponding
\path{dense_ffn_cotenancy_*.json} and
\path{private_transformer_lanes_*.json} files; extraction is in
\path{local_exact_extraction_20260803_172654.json}; the full-basis sweep is in
\path{full_run_r128_log.txt} and
\path{r1024_test_20260622_103122.json}; and the whole-model, pre-MLP, and broken
placement checks are in
\path{archive/smoke/transformer_cotenancy_local_20260812_100546.json}. The
capacity-preserving classifier artifact is
\path{capacity_preserving_wide_20260820_225331.json}. The repository README maps
each artifact to its public runner.

\subsection{Formative co-training at narrow supports}

The formative study applies a regime mask to DistilBERT's final
768-dimensional CLS vector after the encoder. It therefore tests
support-constrained task learning, not intermediate-FFN parameter confinement.
At each $R$, one gated model trains all regimes while the paired separate-model
control trains four sampled regimes. Each statistic compares those same four
regimes across seeds 42, 123, 256, 456, and 789.

\begin{table}[ht]
\centering
\caption{Five-seed paired post-encoder co-training study. Gap is separate minus
gated accuracy. Intervals use Student's $t$ distribution with four degrees of
freedom.}
\label{tab:formative}
\small
\begin{tabular}{@{}rrrrr@{}}
\toprule
$R$ & Dims/regime & Gap (pp) & 95\% CI (pp) & Paired $p$ \\
\midrule
8   & 96 & 0.348 & [$-0.229$, 0.925] & 0.1692 \\
16  & 48 & 0.268 & [$-0.301$, 0.837] & 0.2616 \\
24  & 32 & 0.208 & [$-0.324$, 0.740] & 0.3392 \\
32  & 24 & 0.267 & [$-0.345$, 0.879] & 0.2920 \\
64  & 12 & 0.143 & [$-0.041$, 0.326] & 0.0967 \\
96  & 8  & 0.376 & [0.045, 0.708]    & 0.0345 \\
128 & 6  & 0.272 & [$-0.208$, 0.752] & 0.1902 \\
\bottomrule
\end{tabular}
\end{table}

The observed costs are small despite supports as narrow as six coordinates.
Only $R=96$ crosses an unadjusted 0.05 threshold; none crosses the Bonferroni
family-wise threshold $0.05/7\approx0.0071$. Failure to reject zero is not proof
of equality, and exact equality is not expected between independently optimized
model inscriptions. The defensible result is the measured effect size: in this
formative study, support-constrained co-training remains within 0.38 percentage
points of the sampled separate-model controls at every tested ratio.

{\raggedright
\noindent\textit{Replay.} Use the legacy post-encoder surface in
\path{experiments/reproduce_distilbert_classification.py} with
the reported five-seed rows retained in
\path{experiments/results/series1_results_20260803_202919.json}.\par}

\subsection{Holding per-regime classifier capacity fixed}

The preceding sweep fixes the 768-dimensional post-encoder representation, so
each support narrows as $R$ grows. That is not the only sizing policy. If a task
needs $q$ hidden coordinates per regime, the architecture may instead choose
$d=qR$. Adding a regime then adds a block without reducing any incumbent
regime's allocated width.

We ported and reran the formative wide-classifier protocol with $R=128$,
$q=10$, and therefore $d=1{,}280$. The two-layer MLP has 256 one-hot synthetic
observations and 256 disjoint answers, two per regime. It is co-trained for
2,000 full-batch epochs using a post-activation binary gate. The sparse
selected-support implementation is also compared with explicit dense-mask
execution.

The owning key produces the expected answer for 256/256 observations. Across
all 32,512 pairs of a foreign observation and wrong regime key, the expected
foreign answer is never produced. Activating all 128 regimes simultaneously
produces 249/256 correct answers, and selected-support versus dense-mask logits
have maximum absolute difference 0. The full runner and JSON artifact are
checked in.

This result is deliberately narrow. It is a deterministic synthetic
memorization and scalability test with no held-out set, and the zero wrong-key
hit count is an empirical negative control rather than a privacy theorem. It
does show why ``$R$ regimes'' need not imply ``$1/R$ of a fixed model'' for
modular classifiers: each regime retains the same ten-coordinate allocation
because the governed layer grows with $R$. The cost is linear parameter storage,
not lost per-regime width.

\noindent\textit{Replay.} The self-contained runner is
\path{experiments/capacity_preserving_wide_classifier.py}; the reported run is
\path{experiments/results/capacity_preserving_wide_20260820_225331.json}.

\subsection{Deployment service-consolidation consequence}

An $R=8$ DistilBERT SST-2 deployment benchmark compares eight complete model
services with three single-backbone layouts. Eight services require 1.071 GB
of checkpoint tensors and 1.103 GB of resident CUDA parameter and buffer
state. A frozen shared backbone with eight zero-initialized private adapter
slots retains the original 91.06\% accuracy using 151.6 MB and 156.2 MB,
respectively. This is a deployment control, not evidence that untrained
adapters add utility.

Post-hoc one-eighth FFN slicing reduces accuracy to 67.66\%. Physically
extracting the authorized slices preserves the sliced logits within the
declared fp16 tolerance and raises measured throughput from 1,585 to 2,872
samples/s. Every scored request passes through execution authority and malformed
authority makes zero model calls. Thus the measurement supports consolidating
a frozen backbone plus private modules, while directly rejecting the stronger
claim that naive narrowing of a trained dense Transformer preserves utility.
Timing is one serial stream on one NVIDIA T4 and is not a concurrency result.

\noindent\textit{Replay.} Run
\path{experiments/modal_distilbert_service_consolidation.py}; the full artifact
is listed below.
\par\smallskip\noindent{\footnotesize
\path{experiments/results/distilbert_service_consolidation_20260821T135530_707107Z.json}.}
\par

\subsection{Strict intermediate-FFN cotenancy}

\subsubsection{Protocol}

We evaluate DistilBERT-base~\citep{sanh2019distilbert} on AG
News~\citep{zhang2015agnews}. Attention, embeddings, normalization, residual
parameters, and every other non-governed shared weight are frozen. A binary gate
is placed on the 3,072-dimensional expanded activation immediately before every
FFN down projection. For regime $r$, training may update only the active
first-projection hidden rows, corresponding hidden bias entries, matching
second-projection columns in PyTorch storage convention, and a private
regime-selected classifier. A regime-scoped Adam implementation updates moments
only on authorized entries.

The model is trained on 8,000 examples for two epochs and evaluated on 2,000
held-out examples at seeds 42, 123, and 256. Label permutations supply
regime-specific functional identities. The $R=1$ arm is a matched full-FFN
capacity reference under the same frozen-backbone protocol.

\begin{table}[ht]
\centering
\caption{Strict intermediate-FFN cotenancy. Accuracy is mean $\pm$ sample
standard deviation over three seeds.}
\label{tab:strict}
\small
\begin{tabular}{@{}rrrrr@{}}
\toprule
$R$ & Dims/tenant/layer & Owning accuracy & Drop (pp) & Max unauthorized delta \\
\midrule
1  & 3,072 & $91.28\%\pm0.33\%$ & 0.00 & 0 \\
2  & 1,536 & $90.15\%\pm0.39\%$ & 1.13 & 0 \\
4  &   768 & $88.93\%\pm0.30\%$ & 2.36 & 0 \\
8  &   384 & $87.47\%\pm0.04\%$ & 3.82 & 0 \\
16 &   192 & $86.21\%\pm0.08\%$ & 5.07 & 0 \\
\bottomrule
\end{tabular}
\end{table}

\subsubsection{Interpretation}

The unauthorized-delta column is the maximum over frozen shared parameters,
off-partition FFN parameters, off-partition first and second optimizer moments,
and inactive private classifiers. Every category is bit-exact zero in every
run. The utility cost is empirical and architecture-specific: halving each FFN
retains $90.15\%$ accuracy, 1.13 percentage points below the full-capacity
reference, while the 192-coordinate slice at $R=16$ incurs a 5.07-point drop.

This experiment does not prove parity and does not need to. Independently
optimized model inscriptions have ordinary variance. The experiment measures a
capacity curve under a strict state boundary. The exact result is inactive-state
preservation; the measured result is owning-task utility.

A current-library seed-42 rerun at $R=8$ independently records $87.756\%$
owning accuracy and exact zero for the same frozen, off-partition,
optimizer-moment, and inactive-classifier deltas. Scored evaluation first
passes through the source-visible research preflight's model, dimension, and regime checks;
wrong model, dimension, regime, and empty-authority probes make zero model
calls. This one seed corroborates the boundary but does not replace the
three-seed utility estimate above.

\noindent\textit{Replay.} Run
\path{experiments/modal_dense_ffn_cotenancy.py}; the three-seed artifacts are
the \path{experiments/results/dense_ffn_cotenancy_*.json} files and their
aggregate is \path{experiments/results/transformer_cotenancy_summary.json}.

\subsection{Restoring private capacity}

The strict disjoint-slice construction gives each regime
$d_{\mathrm{ff}}/R$ trainable coordinates. A complementary $R=4$ experiment
asks what happens when each regime instead receives a complete private module
on the same frozen shared DistilBERT backbone. The adapter design adds a
64-dimensional residual adapter after every block; the expert design adds a
complete 3,072-dimensional private FFN after every block. The regime-selected
classifier is private in both designs, while attention remains frozen and
shared.

\begin{table}[ht]
\centering
\caption{Complete private lanes on frozen DistilBERT at $R=4$. Accuracy is mean
$\pm$ sample standard deviation over seeds 42, 123, and 256.}
\label{tab:private-lanes}
\small
\begin{tabular}{@{}lrrrr@{}}
\toprule
Design & Params/regime & Owning acc. & Wrong-key acc. & Max delta \\
\midrule
64-dim residual adapters & 0.59M & $90.47\%\pm0.18\%$ &
$3.18\%\pm0.06\%$ & 0 \\
3,072-dim private experts & 28.33M & $90.09\%\pm0.14\%$ &
$3.30\%\pm0.05\%$ & 0 \\
\bottomrule
\end{tabular}
\end{table}

The adapter is the best measured capacity tradeoff: it recovers most of the
full-FFN utility with much less private state than the expert. The max-delta
column covers the frozen shared backbone and every inactive lane. Wrong-key
accuracy is a functional negative control induced by the regime-specific label
permutations, not the source of the structural claim. This design restores
private capacity rather than dividing one trainable private module by $R$; its
cost is per-regime storage.

A current-library seed-42 rerun records $90.66\%$ adapter and $90.15\%$
expert owning accuracy, with $3.11\%$ and $3.28\%$ wrong-key accuracy.
Shared-backbone and inactive-lane deltas remain exactly zero, and malformed
execution authority makes zero model calls. These one-seed values corroborate,
rather than replace, the historical three-seed estimates.

\noindent\textit{Replay.} Run
\path{experiments/modal_private_transformer_lanes.py}; the reported adapter and
expert runs are \path{experiments/results/private_transformer_lanes_*.json}.

\section{Exact Authorized Extraction}

Let $A_r=\supp(\bg_r)$. For a projection $\bW\in\R^{4\times768}$,
input $\mathbf X\in\R^{256\times768}$, and bias $\mathbf b$, compare
\begin{equation}
  (\mathbf X\odot\bg_r)\bW^\top+\mathbf b
  \quad\text{and}\quad
  \mathbf X[:,A_r]\,\bW[:,A_r]^\top+\mathbf b.
\end{equation}
The extracted expression removes every inactive product rather than computing
it and multiplying by zero.

Across ten seeds, four individual regimes, authorized unions of one through
four regimes, and both 32-bit and 16-bit arithmetic, all 160 comparisons satisfy
their declared numerical criteria. In the canonical 768-dimensional run, all
80 32-bit comparisons are bit-identical. This is an observed kernel-specific
result, not a general floating-point theorem: smaller checks can differ by a
few fp32 units in the last place because the two expressions use different
reduction shapes. Of the 80 16-bit comparisons, 60 are bit-identical and the
maximum absolute difference is 0.015625. Every comparison is within a
conservative forward-error bound computed from the input, weights, support
size, accumulator precision, and bias.

Extraction is both a consequence and a practical use of binary activation. A
runtime may enforce the full mask dynamically, or an authority may export only
the authorized support as a reduced dense artifact. The extracted artifact does
not require the runtime multiplication because excluded coordinates are absent.

\noindent\textit{Replay.} Run \path{experiments/local_exact_extraction.py}; the
reported artifact is
\path{experiments/results/local_exact_extraction_20260803_172654.json}.

\section{Full-Capacity, Zero-Degradation Controls}

The phrase ``zero degradation'' is exact only when the authorized computation
retains all capacity it needs. It should not be inferred from a non-significant
difference between independently trained models. Three completed controls make
the distinction concrete.

\begin{table}[ht]
\centering
\caption{Inference-preservation controls. These retain the complete authorized
computation; they do not claim that dividing an FFN into narrower supports is
cost-free.}
\label{tab:zero-cost-controls}
\small
\begin{tabularx}{\textwidth}{@{}p{0.27\textwidth}X@{}}
\toprule
Control & Result and interpretation \\
\midrule
Atomic whole-model activation & Maximum authorized output difference 0 after one
unchanged model call; five wrong scopes make zero calls \\
Authorized extraction & Same authorized linear computation with inactive products
removed: fp32 80/80 bit-identical; fp16 80/80 within tolerance \\
Keyed full-basis placement & Same trained function in a different coordinate basis:
maximum accuracy gap 0.00 pp at $R=4$--128 and eight sampled placements at
$R=1{,}024$ \\
\bottomrule
\end{tabularx}
\end{table}

The atomic authorization probe binds tenant, sub-id, model digest, and operation
before execution. Its owning path makes one call and is numerically identical to
direct invocation. Random, wrong-tenant, wrong-sub-id, wrong-model, and
wrong-operation credentials make zero model calls. This is Binary Activation at
the complete-model boundary: it preserves model quality because it does not
alter the model.

The extraction result is likewise an implementation-equivalence result. It
removes only products that the authorized binary mask already forces to zero.
It says that a runtime may materialize the active support without changing the
authorized linear result; it does not compare a narrow support with a wider
model.

Finally, a full-capacity DistilBERT control conjugates every residual-facing
parameter consistently into a keyed permutation basis. The base model and all
tested placements score $90.25\%$ on the same 7,600-example test set, with
exactly zero prediction-level accuracy gap for all placements at $R=4$ through
$R=128$ and eight sampled placements at $R=1{,}024$. No dimension is removed in
this control, so the result is consistent with exact function preservation.
Conjugation alone is not Binary Activation: it creates neither distinct tenant
knowledge nor a zero on non-owning coordinates. It is included only to show
that keyed placement itself has zero inference cost when the full basis remains
available. Tenant confinement still requires a correctly placed gate and the
capacity tradeoffs measured above.

\noindent\textit{Replay.} Atomic whole-model authorization is exercised by
\path{experiments/local_transformer_cotenancy_suite.py}; extraction by
\path{experiments/local_exact_extraction.py}; and the keyed full-basis control
by \path{experiments/orthogonal_superposition_sweep.py}. Retained outputs are
listed in \path{experiments/results/README.md}.

\section{Placement and Bypass Negative Controls}

A deterministic local protocol run tests both the valid gate and two invalid
constructions. At a correctly placed disjoint gate, the maximum cross-mask value
is exactly zero. Inserting LayerNorm between the source assignment and the
target gate produces a nonzero target coordinate of 0.12598. Adding an ungated
residual route carries the excluded source with maximum magnitude 3.0. The
controls therefore fail in the direction predicted by the threat model.

A separate pre-MLP probe gates a 32-dimensional FFN input before a 32-to-64
first projection at $R=4$. The full gated and physically extracted projections
have maximum difference 0; inactive input gradients and aligned inactive
first-projection-row gradients are also exactly zero. This is evidence for the
declared pre-MLP surface only, not for hidden/down-projection partitioning.

These controls matter because an always-zero test could otherwise reflect a
dead harness. The observed nonzero failures show that the implementation detects
the normalization and bypass conditions that invalidate a broader isolation
claim.

\noindent\textit{Replay.} Run
\path{experiments/local_transformer_cotenancy_suite.py}. The retained protocol
artifact is
\path{experiments/results/archive/smoke/transformer_cotenancy_local_20260812_100546.json}.

\section{Threat Model and Limits}

The in-scope adversary is a caller that may submit chosen inputs, possess its
own credential, alter an asserted regime identifier, or guess coordinate
assignments. The trusted runtime verifies the credential, resolves the regime,
obtains the authority-selected mask, and applies the gate at the declared FFN
surface. It does not accept a caller-selected mask.

The following remain trusted or separately controlled:
\begin{itemize}
  \item credential issuance, verification, and revocation;
  \item assignment-key custody and credential-to-mask dispatch;
  \item model and tensor-axis integrity;
  \item gate placement and closure of residual or other bypass routes;
  \item optimizer state and weight-decay behavior; and
  \item caches, logs, backups, retrieval stores, adapters, and other serving
  state outside the declared FFN parameter surface.
\end{itemize}

A wrong credential cannot activate the owning regime's coordinates under
correct dispatch because disjointness supplies an exact zero at that gate. That
does not imply that the complete output is zero: the caller's authorized
coordinates remain active. Nor does it prove statistical independence of shared
representations. Attention and other shared components may compute common
features before or after the gated FFN; their parameters are not represented as
private by this theorem.

Improper gate placement can nullify the intended isolation. If a protected
signal is mixed into active coordinates before masking, or if an ungated branch
carries it around the gate, multiplication still zeros the immediate tensor but
does not isolate the complete computation. A conforming implementation must
state exactly which activation, parameter slices, optimizer state, and runtime
state it claims to confine.

\section{Related Work}

Task-specific masking methods such as PackNet~\citep{mallya2018packnet},
Piggyback~\citep{mallya2018piggyback}, and Supermasks in
Superposition~\citep{wortsman2020supermasks} demonstrate that subnetworks can
share a parameterized model. Binary activation differs in joining a
credential-derived disjoint assignment to both inference and training-time
support. LoRA~\citep{hu2022lora} provides efficient private trainable state, but
does not by itself supply an identity-bound binary activation boundary.

Formal verification of neural systems commonly models robustness or system
properties at higher levels~\citep{seshia2018formal}. Here the principal formal
results are deliberately small: exact consequences of a binary Hadamard gate at
a declared FFN boundary.

\section{Conclusion}

Binary activation connects identity to neural computation. A verified
credential selects a disjoint FFN support; the Hadamard gate determines which
coordinates can contribute to inference and where loss gradients can flow. At
the correctly placed intermediate FFN activation, the local exclusion is exact.
With aligned optimizer state, the corresponding parameter updates are exact as
well.

The experiment set separates effects that a single accuracy table cannot.
Formative co-training uses supports as narrow as six coordinates with observed
costs below 0.38 percentage points. The wide-classifier experiment instead
holds per-regime width fixed while scaling $d=qR$, and recovers 256/256 owning
predictions at $R=128$ on its synthetic memorization task. Strict
intermediate-FFN cotenancy keeps
inactive state bit-exact unchanged while utility declines as each regime
receives less FFN capacity. Complete private lanes recover utility by restoring
per-regime capacity at a storage cost. Exact extraction and full-capacity
controls then show where zero inference degradation is the right statement:
when authorization or reparameterization leaves the complete authorized
computation available. Deliberately broken placement controls confirm that
normalization and ungated residual paths invalidate broader claims.

These results do not require complete Transformer partitioning, generative
validation, or a statistical assertion of zero leakage. They require a
correctly placed binary gate, aligned optimizer state, and a trusted
credential-to-mask runtime. Capacity-preserving controls are reported as such,
not as evidence that every divided FFN partition is free.

In capability-infrastructure terms, the contribution is an identity-addressable
internal support to which a scoped grant can refer; the runtime and lockbox
architecture that issues and redeems those grants is treated separately.

Code, artifacts, and Lean proofs are available at
\url{https://github.com/sekosai/schemen-gate/tree/main/research/cdp}.

\small
\bibliographystyle{plainnat}
\begin{thebibliography}{8}

\bibitem[Hu et~al.(2022)]{hu2022lora}
E.~J. Hu, Y.~Shen, P.~Wallis, et~al.
\newblock LoRA: Low-rank adaptation of large language models.
\newblock In \emph{ICLR}, 2022.
\newblock \url{https://arxiv.org/abs/2106.09685}

\bibitem[Mallya and Lazebnik(2018)]{mallya2018packnet}
A.~Mallya and S.~Lazebnik.
\newblock PackNet: Adding multiple tasks to a single network by iterative pruning.
\newblock In \emph{CVPR}, 2018.
\newblock \url{https://arxiv.org/abs/1711.05769}

\bibitem[Mallya et~al.(2018)]{mallya2018piggyback}
A.~Mallya, D.~Davis, and S.~Lazebnik.
\newblock Piggyback: Adapting a single network to multiple tasks by learning to mask weights.
\newblock In \emph{ECCV}, 2018.
\newblock \url{https://arxiv.org/abs/1801.06519}

\bibitem[Sanh et~al.(2019)]{sanh2019distilbert}
V.~Sanh, L.~Debut, J.~Chaumond, and T.~Wolf.
\newblock DistilBERT, a distilled version of BERT: smaller, faster, cheaper and lighter.
\newblock \emph{arXiv:1910.01108}, 2019.
\newblock \url{https://arxiv.org/abs/1910.01108}

\bibitem[Seshia et~al.(2018)]{seshia2018formal}
S.~A. Seshia, A.~Desai, T.~Dreossi, et~al.
\newblock Formal specification for deep neural networks.
\newblock In \emph{ATVA}, 2018.
\newblock \url{https://arxiv.org/abs/1710.01688}

\bibitem[Wortsman et~al.(2020)]{wortsman2020supermasks}
M.~Wortsman, V.~Ramanujan, R.~Liu, et~al.
\newblock Supermasks in superposition.
\newblock In \emph{NeurIPS}, 2020.
\newblock \url{https://arxiv.org/abs/2006.14769}

\bibitem[Zhang et~al.(2015)]{zhang2015agnews}
X.~Zhang, J.~Zhao, and Y.~LeCun.
\newblock Character-level convolutional networks for text classification.
\newblock In \emph{NeurIPS}, 2015.
\newblock \url{https://arxiv.org/abs/1509.01626}

\end{thebibliography}
\end{document}
